\section{Considerações Finais}
\label{sec:consideracoes-finais}

O trabalho entrega requisitos priorizados, três objetivos GQM, sete questões, quatorze métricas e procedimentos reproduzíveis para comparar a consulta de infrações nos dois canais. As segundas leituras elegíveis das sessões de 21 e 22/11 foram concluídas e os níveis calculados. G1 não foi atendido em nenhum canal: a lista apresentou 5/15 campos e a ausência de R14 levou a 0/1 infração integralmente paritária. G2 permanece inconclusivo porque os instrumentos de volume não são equivalentes, embora os tempos medidos tenham nível Excelente e o portal tenha sido mais rápido na amostra V4. G3 foi atendido para o portal nos contextos medidos e não atendido para o aplicativo no escopo observado, devido à falha 0/3 em CT05 com a versão 2.6.3. A condição degradada no aplicativo e o CT04 não foram avaliados; não se presume resultado para essas lacunas nem se atribui causa às diferenças entre versões e aparelhos.

Foi entregue uma proposta de processo para diagnóstico e regressão do bloqueio no aplicativo, com ficha, matriz de testes e critérios de reavaliação, ligada às decisões institucionais D1 e D4. Sua entrega não comprova correção do serviço. Os consentimentos foram aplicados, conforme confirmação de Hugo em 22/11, com vias assinadas privadas. As conclusões permanecem restritas aos veículos, dispositivos, versões e condições observados; não representam todos os cidadãos nem certificam o serviço do Detran-DF.
